\documentclass[12pt]{article}
\usepackage[utf8]{inputenc}
\usepackage[T1]{fontenc}
\usepackage[margin=1in]{geometry}
\usepackage{setspace}
\usepackage{xcolor}
\usepackage[colorlinks=true, urlcolor=blue!60!black, linkcolor=blue!60!black]{hyperref}
\setlength{\parindent}{2em}
\setlength{\parskip}{1.25em}
\renewcommand{\baselinestretch}{1.0}
\begin{document}
\begin{center}
{\LARGE \textit{Martin v. Hunter's Lessee}, 14 U.S. (1 Wheat.) 304 (1816)\par}
\vspace{0.8em}
{\large Case Brief --- Four-Page Edition\par}
\end{center}
\vspace{1em}

\section{Memory Jogger}

The Supreme Court held that Article~III and Section~25 of the Judiciary Act of 1789 authorize the Supreme Court to review final state-court civil judgments involving federal constitutional, statutory, or treaty questions, to preserve national supremacy and uniformity of federal law (\href{https://supreme.justia.com/cases/federal/us/14/304/}{Justia}).

\section{Detailed Case Facts}

Lord Fairfax held title under royal letters patent to the Northern Neck of Virginia. When he died in 1781, he devised the lands to his nephew Denny Martin (later Denny Fairfax), a British subject who lived in England throughout the Revolution. Virginia enacted laws restricting alien landownership and confiscating Loyalist estates, though the parcel at issue had never been formally escheated. In 1789 Virginia granted the parcel to David Hunter. The 1783 Treaty of Peace and the 1794 Jay Treaty protected British subjects' landholdings. Denny Fairfax's claim passed to Philip Martin. Chief Justice Marshall and relatives had contracted to buy parts of the Fairfax estate, so Marshall recused.

\section{Procedural History}

Hunter's lessee sued in ejectment in 1791; the Winchester District Court ruled for Fairfax in 1794. The Virginia Court of Appeals reversed in 1810. In \textit{Fairfax's Devisee v. Hunter's Lessee} (1813), the Supreme Court, per Justice Story, reversed under Section~25, holding that the Jay Treaty protected Fairfax's title. Led by Judges Spencer Roane and William Cabell, the Virginia Court of Appeals refused to obey the mandate, declaring Section~25 unconstitutional insofar as it reached state courts. On a second writ of error, the Supreme Court reversed on March 20, 1816, and directly affirmed the Winchester judgment.

\section{Judicial Votes}

Reversed, \textbf{6--0}.
\begin{itemize}
\item \textbf{Majority (5):} Justice Story, joined by Justices Washington, Livingston, Todd, and Duvall.
\item \textbf{Concurrence:} Justice Johnson, in the judgment.
\item \textbf{Recused:} Chief Justice Marshall. No dissents.
\end{itemize}

\section{Holding}

Under Article~III and the Supremacy Clause, the Supreme Court's appellate jurisdiction extends to final state-court judgments in civil cases involving federal questions, because Article~III attaches jurisdiction to the \textit{case}, not the \textit{court}. Section~25 is constitutional; a state court's refusal to obey a mandate is itself a reviewable final judgment; and the Court may examine an intertwined state-law title question so that state courts cannot evade review of a treaty claim. Story's argument that Congress must vest the entire judicial power over ``all'' first-tier cases was dicta, and the Court reserved whether it could issue mandamus to a state court.

\section{Analysis of Opinions}

\textbf{Majority (Story).} The Constitution was ordained by ``the people of the United States,'' who could subordinate state powers to national authority, and its general grants must be reasonably construed. Story argued that ``shall be vested'' and ``shall extend'' are imperative, and that ``all'' in the first-tier heads signals mandatory federal jurisdiction; independently, Article~III's appellate grant contains no limit on the tribunal from which an appeal comes. Four arguments reinforced the text: the Supremacy Clause contemplates state judges deciding federal questions, which ``all'' cases can reach only on appeal; federal review of state judicial acts is no greater intrusion than review of state statutes; state attachments may sway state courts, and only one revising tribunal can harmonize ``jarring and discordant judgments''; and the First Congress and long state acquiescence confirmed Section~25.

\textbf{Johnson (concurring).} Viewing the Constitution as a ``tripartite'' compact, Johnson rejected Story's mandatory-vesting and two-tier theories: Congress could leave federal questions to state courts subject to revision. Stressing comity, he said the Court acts ``over persons and cases,'' not over state courts as subordinates, and read Section~25's power to award execution as a way to avoid mandamus to state judges.

\section{Examples: Future Applications}

\begin{enumerate}
\item \textbf{Same-Side: State Criminal Conviction.} A state conviction under a lottery law rejects a defense based on a federal statute. Jurisdiction attaches to the federal question in the case, as \textit{Cohens v. Virginia} (1821) held.
\item \textbf{Same-Side: Defiance of a Mandate.} A state supreme court, on remand, declares a federal preemption statute unconstitutional and refuses judgment. The defiant order is reviewable, and the Court may reverse without mandamus.
\item \textbf{Opposite-Side: Adequate and Independent State Ground.} A state court rules for a landlord both on federal law and on a neutral state waiver rule. The state ground independently supports the judgment, so the Court lacks jurisdiction (\textit{Staub v. City of Baxley}).
\item \textbf{Opposite-Side: State Review of Federal Judgments.} A litigant asks a state court to set aside a federal judgment. \textit{Martin} runs one way; state courts have no reciprocal power (\textit{Ableman v. Booth}).
\item \textbf{Fence-Sitter: Jurisdiction-Stripping.} Congress strips all federal courts of jurisdiction over challenges to state abortion or religion laws. \textit{Martin}'s narrow holding and Johnson's concurrence suggest broad Exceptions Clause power, but Story's mandatory-vesting dicta, championed by Akhil Amar, suggest a violation.
\end{enumerate}

\section{Critique}

\textbf{Originalist and textualist views.} Roane and Cabell argued that Article~III never mentions state courts and that writs of error ``commandeered'' sovereign state judiciaries. Modern originalists vindicate the core holding through the Supremacy Clause, \textit{Federalist No.~82}, and the First Congress. Amar embraces Story's two-tier reading; Martin Redish and Daniel Meltzer say it clashes with the Madisonian Compromise; James Pfander offers an alternative grounded in Congress's power to create inferior tribunals.

\textbf{Progressive and institutional views.} Scholars celebrate \textit{Martin} as the linchpin of national rights, without which federal law could be nullified state by state. Its re-examination of state law was later cabined by the adequate-and-independent-state-ground doctrine (\textit{Murdock v. Memphis}), and critics say it helped build a judicial monopoly on constitutional interpretation.

\section{Key Quotations}

\begin{quote}
``That motive is the importance, and even necessity, of uniformity of decisions throughout the whole United States upon all subjects within the purview of the Constitution. Judges of equal learning and integrity in different States might differently interpret a statute or a treaty of the United States, or even the Constitution itself; if there were no revising authority to control these jarring and discordant judgments and harmonize them into uniformity, the laws, the treaties, and the Constitution of the United States would be different in different States, and might perhaps never have precisely the same construction, obligation, or efficacy in any two States.''
\end{quote}
\noindent --- Justice Story, Opinion of the Court, 14 U.S. (1 Wheat.) at 347--48.

\begin{quote}
``The Constitution of the United States was ordained and established not by the States in their sovereign capacities, but emphatically, as the preamble of the Constitution declares, by `the people of the United States.' ''
\end{quote}
\noindent --- Justice Story, Opinion of the Court, 14 U.S. (1 Wheat.) at 324.

\begin{quote}
``It is the case, then, and not the court, that gives the jurisdiction. If the judicial power extends to the case, it will be in vain to search in the letter of the Constitution for any qualification as to the tribunal where it depends.''
\end{quote}
\noindent --- Justice Story, Opinion of the Court, 14 U.S. (1 Wheat.) at 338.

\end{document}
